% =============================================================================
% Discussion + Limitations -- rewritten 2026-07-27 against the final Results.
% Interpretation only: no quantity appears here that Results has not reported,
% except the two cross-cutting comparisons flagged in the comments below.
% =============================================================================

\section{Discussion}

\paragraph{One localization, two failures.}
Locating the failure in the belief to action mapping does not make it one
phenomenon. Under Bayesian inference the belief is coupled to the action and
underweighted: amplifying it along its own direction recovers a substantial
share of errors, so the pathway exists and carries too little weight. Under
game theoretic structure the same intervention is barely distinguishable from
a matched random perturbation, and the action is far less sensitive to a
belief the model demonstrably holds. These are a gain problem and a
connectivity problem. They share a name and not a mechanism, and only the
first is the kind of defect a scaling or reweighting fix could plausibly
address.

\paragraph{Reactivity degrades the belief; game structure breaks the mapping.}
The two stages are governed by different properties of a scenario, and one
scenario separates them. \textit{Chicken} has a payoff matrix and an opponent,
but its opponent is a fixed yield rate rather than a strategy that responds to
the agent. Its belief forms like a nonstrategic inference problem and its
action collapses like a game. The comparison is direct: \textit{Chicken} and
the reactive scenarios reach almost the same action agreement, \(66.1\) and
\(65.2\), from very different belief levels, \(98.5\) and \(90.2\). What
depresses belief formation is the need to condition on the agent's own
previous action; what breaks the mapping is the strategic framing. Neither
explains the other, and a suite without a nonreactive game could not have told
them apart.

This also bounds the obvious alternative account. From the Bayesian inference
scenarios to the reactive ones the belief falls by roughly nine points while
action agreement falls by more than twice that, and the intervention shows the
action does not track the belief in the game theoretic scenarios at any level.
A weaker belief is present. It is not sufficient.

\paragraph{Both interventions fail on the same decisions.}
Steering and chain of thought act through different routes. Steering operates
on the model's own belief to action pathway; reasoning bypasses that pathway
and substitutes an external, serialized computation. They therefore succeed in
different places, and reasoning recovers errors in exactly the scenarios where
steering cannot.

Where they agree is more informative. Both fail when the model is confidently
wrong. Steering does not move the largest models on their errors, because
those errors sit five to eleven nats from the decision boundary, beyond what
the intervention can cross. Reasoning shows the same profile: its correction
rate falls from \(95.5\%\) on the errors closest to the boundary to \(30.4\%\)
on the most confident quartile for \model{Gemma-4} 31B. Two interventions with
different mechanisms fail on the same decisions, which suggests the binding
constraint is not the route taken but how far the decision already is from
reversing.

\paragraph{What this implies for training.}
The belief is present, linearly decodable at every scale we test, and
therefore already available as a supervised target. Nothing is gained by
teaching these models the structure of the scenarios; that information is
there. What must be trained is a utility sensitive mapping from a represented
belief to an action. Neither more reasoning supervision nor a stronger belief
addresses that, and the degradation of \model{Gemma-2} 2B under chain of
thought is the direct evidence: layering reasoning on top of a broken mapping
does not repair it, it amplifies whatever the mapping already does.

\paragraph{Why localization was necessary.}
An outcome score cannot separate the failures this paper distinguishes. A
model that never formed the belief, a model that formed it and did not act on
it, and a model that acted correctly and then reasoned itself out of it all
produce one wrong answer. On this suite they occur in different models and
different scenarios, and they call for different interventions. Prompting for
more careful inference targets a stage that is not failing. The measurements,
not the outcomes, identify where to intervene.

\subsection{Limitations}

\textbf{\textit{Committed opponents.}}
Opponents follow fixed strategies rather than best responding, so the game
theoretic scenarios are Bayesian decision problems against committed types
rather than equilibrium computations. This is what makes an exact optimal
action available at every step, and it scopes our claims to best response
against a predicted opponent.

\textbf{\textit{The belief evidence is correlational; only steering is causal.}}
Linear decodability establishes that the information is present, not that the
model uses it. The steering intervention supplies the causal evidence, and it
is narrow: one layer, the final prompt position, a single direction per
scenario. Orthogonal and shuffled target directions, and direct verification
that the intervention moves the decoded posterior, remain to be run.

\textbf{\textit{A simple statistic recovers the posterior in the reactive
scenarios.}}
Our prompts print a running tally. In the reactive scenarios a quadratic
function of the joint tally of the agent's previous action and the opponent's
move predicts the posterior better than the residual stream does. The
like-for-like linear comparison favors the representation throughout, but we
do not claim the network computes a quantity no simple function of the tally
could.

\textbf{\textit{A scenario we designed out.}}
An earlier version of the suite included the Iterated Prisoner's Dilemma. We
excluded it after establishing that its optimal action is not belief driven:
defection strictly dominates in the stage game under the payoffs used, so
restoring belief relevance requires a finite horizon continuation term, which
makes the optimal action a function of the remaining horizon as well as the
posterior. A predictor that sees only the step index and never the history
reproduces the resulting labels on \(70.5\%\) of steps against a base rate of
\(56.3\%\).

\textbf{\textit{Scope.}}
All scenarios use a binary action space and a fixed horizon of ten steps.
Histories are generated by a fixed behavior policy, so our claims concern
observer style decision making: whether a model responds rationally to
provided evidence, not whether its own policy would have generated
informative evidence.
